En el contexto de la resolución de un problema de programación  lineal, $P$, mediante el método de las dos fases, explicar justificadamente:
\begin{enumerate}
    \item ¿Cuál es la interpretación de la función objetivo correspondiente al problema de la primera fase, $P'$ en relación con el problema $P$?
    \item ¿Qué interpretación tiene el hecho de que la solución óptima de $P'$ tenga función objetivo diferente de cero?
\end{enumerate}